% ============================================================
% Sección 1: Introducción
% ============================================================

\section{Introducción}
\label{sec:intro}

El speckle óptico es el patrón granular de alta intensidad que emerge cuando luz coherente
---como la de un láser--- se dispersa sobre una superficie rugosa.
Este fenómeno está presente en aplicaciones de metrología óptica, tomografía de coherencia
óptica, imagenología biomédica, criptografía óptica y comunicaciones de fibra
\citep{Goodman2007_Speckle}.
Su predicción precisa requiere resolver la ecuación de Helmholtz con condiciones de frontera
estocásticas, un problema que los métodos numéricos clásicos abordan a un costo computacional
prohibitivo.

Los métodos basados en malla, como el Método de Elementos Finitos (FEM), son precisos pero
imponen restricciones severas: la malla debe tener resolución espacial inferior a la longitud
de onda del láser para capturar los detalles finos del speckle \citep{Jin2014_FEM}.
Esto conduce a sistemas con millones de incógnitas por realización, y cada nuevo perfil de
rugosidad superficial exige resolver el sistema completo desde cero.
Para estudios estadísticos que requieren miles de realizaciones, el costo se vuelve intratable.

Las Redes Neuronales Informadas por Física (\textit{Physics-Informed Neural Networks},
PINNs), introducidas por \citet{Raissi2019_PINNs}, proponen un paradigma alternativo
\textit{libre de malla} (\textit{mesh-free}): incorporar las leyes físicas directamente en la
función de pérdida de la red, forzando simultáneamente el ajuste a las condiciones de frontera
y el cumplimiento de la ecuación diferencial en el dominio.
Una vez entrenado, el modelo evalúa el campo completo en una malla de
$100 \times 100$ puntos en $\approx 1.3\,\mathrm{ms}$ (GPU NVIDIA RTX~5050),
actuando como un sustituto computacionalmente eficiente del solver numérico.
La eficacia de este paradigma ha sido demostrada en problemas de dispersión de ondas
\citep{Fang2020_scattering}, nano-óptica e inversión de parámetros
\citep{Chen2020_nanoptics}, y el marco de referencia de PINNs ha sido consolidado
en revisiones de \citet{Karniadakis2021_review} y \citet{Cuomo2022_review}.

Un obstáculo crítico de las PINNs estándar para la ecuación de Helmholtz es la inestabilidad
de las derivadas de segundo orden con activaciones tipo ReLU o tanh.
Las \textit{Sinusoidal Representation Networks} (SIREN), propuestas por
\citet{Sitzmann2020_SIREN}, resuelven este problema mediante activaciones
$\phi(z) = \sin(\omega_0 z)$ con una inicialización que preserva la distribución de
activaciones en toda la red, convirtiendo la arquitectura en la elección natural para
ecuaciones de onda.
Un estudio reciente de \citet{Schoder2024_PINN_acoustics} validó la viabilidad de las PINNs
para la ecuación de Helmholtz 3D, reportando un error $L^2$ de $2.5\%$ respecto a soluciones
FEM de referencia.
Sin embargo, hasta donde alcanza nuestro conocimiento, la aplicación sistemática de este
paradigma a la simulación del speckle óptico de alta resolución no ha sido abordada.

En este trabajo presentamos una implementación PINN-SIREN para la simulación acelerada del
speckle óptico que cubre este vacío.
La contribución central es la resolución de la ecuación de Helmholtz bidimensional con campo
complejo y condición de frontera de fase aleatoria $\phi(x) \sim \mathcal{U}(0, 2\pi)$,
que modela una superficie con rugosidad $\sigma_h \gg \lambda$.
Demostramos que la misma arquitectura SIREN $5\times128$ que resuelve la propagación de onda
plana (Error $L^2 = 0.171\%$) es capaz de generar patrones de speckle estadísticamente
válidos según el criterio de \citet{Goodman2007_Speckle}: contraste
$C = \sigma_I / \langle I \rangle \approx 1$ y distribución de intensidades exponencial
negativa, confirmando speckle completamente desarrollado.

El artículo está organizado como sigue.
La \cref{sec:teoria} presenta el marco teórico: la ecuación de Helmholtz, la formulación PINN
y la arquitectura SIREN, y la estadística del speckle.
La \cref{sec:metodologia} describe la formulación del problema, el muestreo LHS y la
estrategia de optimización Adam + L-BFGS.
La \cref{sec:resultados} reporta los resultados en tres etapas: validación 1D, validación 2D
y generación de speckle.
La \cref{sec:conclusiones} sintetiza las contribuciones y señala direcciones futuras.
